\section{Skills 与 MCP/Tools 的正交关系}

\subsection{两种不同的 Context 位置}

Skills 和 MCP/Tools 在技术架构上是\textbf{正交独立}的，它们占据 API 请求中完全不同的位置：

\begin{importantbox}{Skills vs. Tools 的位置差异}
\begin{itemize}
    \item \textbf{Skills}：出现在 \texttt{message} 的 \texttt{content} 字段中，具体是注入的第三条 User 消息，与 agents.md 放在一起
    \item \textbf{MCP / 内置 Tools}：出现在 API 的 \texttt{tools} 字段中，以 Tools Schema 的形式暴露给模型
\end{itemize}
两者处于完全不同的位置，\textbf{不存在} Skills 出现后简化了 MCP/Tools 所占 Context 的情况（除非显式设计）。
\end{importantbox}

\subsection{Skills 的渐进式加载回顾}

以 Codex 为例，Skills 的加载流程：

\begin{enumerate}
    \item 预注入的 User 消息中包含 Skills 的索引——每个 Skill 的名称和简短描述
    \item 用户发出第一个 Query 后，远端模型根据 Query 匹配最合适的 Skill
    \item 如果决定使用某个 Skill，模型通过 Function Call 执行一个 Shell 命令（如 \texttt{cat}），将该 Skill 对应的 skills.md 完整加载
    \item 完整的 Skill 内容作为 Function Output 消息加入 Context
\end{enumerate}

\begin{knowledgebox}{Skills 可以引用 Tools}
Skills 的文档中可以提及具体的 Tools——说明什么时候应该使用什么工具。但这只是指引性的文本，实际的 Tool 定义仍然需要通过 API 的 \texttt{tools} 字段以 Schema 形式暴露给模型。Skills 是"说明书"，Tools 是"工具箱"，两者配合使用。
\end{knowledgebox}

\subsection{本章小结}

Skills 和 MCP/Tools 在 API 层面占据不同位置——Skills 在消息内容中，Tools 在工具字段中。两者技术上正交独立，但在实际使用中协同配合。


